% !TEX program = xelatex
\documentclass[11pt,a4paper]{article}

\usepackage{fontspec}
\usepackage{polyglossia}
\usepackage{amsmath,amssymb,amsfonts}
\usepackage{xcolor}
\usepackage{graphicx}
\usepackage{geometry}

\geometry{left=2cm,right=2cm,top=2cm,bottom=2cm}

\setdefaultlanguage{persian}
\setotherlanguage{english}
\setmainfont{Vazir}

\newcommand{\en}[1]{\textenglish{#1}}
\newcommand{\separator}{\vspace{1em}\noindent\rule{\linewidth}{0.4pt}\vspace{1em}}
\newcommand{\thickseparator}{\vspace{1em}\noindent\rule{\linewidth}{1pt}\vspace{1em}}

\title{پاسخ تمرین بهینه‌سازی ترکیبیاتی: مسأله تخصیص خطی}
\author{محمد باغبان زاده \\ شماره دانشجویی: شناسهٔ دانشجو}
\date{\today}

\linespread{1.35}

\begin{document}

\maketitle
\thickseparator

\section{تعریف مسأله تخصیص خطی (\en{Linear Assignment Problem})}

مسأله تخصیص خطی یکی از مشهورترین مسائل بهینه‌سازی ترکیبیاتی است. فرض کنید $n$ کار و $n$ ماشین (یا فرد) داریم. هزینه (یا سود) اختصاص کار $i$ به ماشین $j$ برابر با $c_{ij}$ است. هدف آن است که هر کار دقیقاً به یک ماشین و هر ماشین دقیقاً یک کار اختصاص یابد، به گونه‌ای که مجموع هزینه‌های تخصیص کمینه (یا بیشینه) گردد.

به بیان ریاضی، مجموعه $X$ از تخصیص‌های یک‌به‌یک میان دو مجموعه $I=\{1,\dots,n\}$ و $J=\{1,\dots,n\}$ به صورت جایگشت $\pi$ از $\{1,\dots,n\}$ مدل می‌شود. مسأله عبارت است از:
\[
\min_{\pi \in S_n} \sum_{i=1}^{n} c_{i,\pi(i)},
\]
که $S_n$ مجموعه همه جایگشت‌های $n$ عنصری است.

\separator
\newpage
\section{مدل‌سازی به صورت برنامه‌ریزی عدد صحیح}

برای هر زوج $(i,j)$ متغیر دودویی $x_{ij}$ تعریف می‌کنیم:
\[
x_{ij} = \begin{cases}
1 & \text{اگر کار $i$ به ماشین $j$ اختصاص یابد}\\
0 & \text{در غیر این صورت}
\end{cases}
\]

مدل برنامه‌ریزی عدد صحیح $0$-$1$ به شکل زیر است:
\begin{align}
\min \quad & \sum_{i=1}^{n}\sum_{j=1}^{n} c_{ij} x_{ij} \label{eq:obj}\\
\text{s.t.} \quad & \sum_{j=1}^{n} x_{ij} = 1, \qquad \forall i=1,\dots,n \label{eq:assign1}\\
& \sum_{i=1}^{n} x_{ij} = 1, \qquad \forall j=1,\dots,n \label{eq:assign2}\\
& x_{ij} \in \{0,1\}, \qquad \forall i,j \label{eq:binary}
\end{align}
محدودیت \eqref{eq:assign1} تضمین می‌کند هر کار دقیقاً یک ماشین دریافت کند و محدودیت \eqref{eq:assign2} تضمین می‌کند هر ماشین دقیقاً یک کار انجام دهد.

\separator
\newpage
\section{پیچیدگی: آیا مسأله تخصیص خطی \en{NP-hard} است؟}

خیر، مسأله تخصیص خطی \textbf{\en{NP-hard} نیست} و در زمان چندجمله‌ای قابل حل است. دلایل اصلی به شرح زیر است:

\begin{enumerate}
\item \textbf{وجود الگوریتم‌های چندجمله‌ای}: مشهورترین الگوریتم، الگوریتم مجارستانی (\en{Hungarian algorithm}) است که این مسأله را در زمان $O(n^3)$ حل می‌کند. همچنین می‌توان آن را به عنوان یک نمونه خاص از مسأله کمترین هزینه جریان (\en{Min-Cost Flow}) با ظرفیت‌های $1$ و $n$ گره مبدأ و $n$ گره مقصد در نظر گرفت که برای آن نیز الگوریتم‌های چندجمله‌ای وجود دارد.

\item \textbf{یکپارچگی (\en{Integrality}) ماتریس قیود}: ماتریس ضرایب قیود \eqref{eq:assign1} و \eqref{eq:assign2} یک ماتریس \textbf{تمام‌اً یک‌پارچه} (\en{Totally Unimodular}) است. یکی از خواص مهم ماتریس‌های تمام‌اً یک‌پارچه آن است که رأس‌های چندوجهی $\{x \ge 0 : Ax = b\}$ برای بردار صحیح $b$، دارای مختصات صحیح هستند. در نتیجه، \textbf{برنامه‌ریزی خطی حاصل از relaxing} کردن شرط صحیح بودن متغیرها ($0 \le x_{ij} \le 1$) یک جواب بهینه با مقادیر $0$ یا $1$ تولید می‌کند. بدین ترتیب، مسأله تخصیص خطی را می‌توان با حل یک برنامه‌ریزی خطی (که در زمان چندجمله‌ای قابل حل است) دقیقاً حل کرد، بدون این‌که نیاز به حل یک \en{ILP} عمومی باشد.

\item \textbf{جایگاه در کلاس $P$}: وجود الگوریتم چندجمله‌ای (چه ترکیبیاتی مانند مجارستانی و چه مبتنی بر \en{LP}) ثابت می‌کند که مسأله تخصیص خطی به کلاس پیچیدگی $P$ تعلق دارد و نمی‌تواند \en{NP-hard} باشد، مگر آن‌که $P = NP$.
\end{enumerate}

بنابراین، مسأله تخصیص یک مثال برجسته از مسائلی است که \en{IP} آن با وجود متغیرهای دودویی، ذاتاً «آسان» است و سخت‌ی \en{NP-hardness} عمومی برنامه‌ریزی عدد صحیح را به ارث نمی‌برد.

\separator
\newpage
\section{پرسش تحلیلی: رابطه یک نمونه حل‌شدنی چندجمله‌ای با \en{NP-hardness} عمومی}

\textbf{اگر یک مسأله \en{integer programming} داشته باشیم که جوابش در زمان چندجمله‌ای به‌دست می‌آید، آیا می‌توان نتیجه گرفت که برنامه‌ریزی عدد صحیح \en{NP-hard} نیست؟}

پاسخ: \textbf{خیر، این نتیجه‌گیری نادرست است.} 

برنامه‌ریزی عدد صحیح عمومی ($IP$) یک چارچوب مدل‌سازی بسیار گسترده است. \en{NP-hard} بودن آن به این معنا است که \textbf{بدترین حالت} مسأله (یعنی وجود حداقل یک خانواده از نمونه‌ها) نیازمند زمان نمایی است، مگر $P=NP$. این خاصیت با ارائه یک نمونه خاص که بتوان آن را در زمان چندجمله‌ای حل کرد، نقض نمی‌شود. به عبارت دیگر:

\begin{itemize}
\item \textbf{\en{NP-hardness} یک ویژگی کل کلاس مسائل است، نه یک نمونه خاص.} یک مسأله \en{NP-hard} می‌تواند زیرمسائل ساده (قابل حل در $P$) فراوانی داشته باشد. برای مثال:
\begin{itemize}
\item مسأله کوله‌پشتی صفر و یک \en{NP-hard} است، اما اگر وزن همه اشیا یکسان باشد با یک مرتب‌سازی ساده و انتخاب حریصانه حل می‌شود.
\item مسأله پوشش رأسی \en{NP-hard} است، اما برای گراف‌های دوبخشی در زمان چندجمله‌ای (از طریق تطابق ماکزیمم) حل می‌شود.
\item مسأله تخصیص خطی که در بالا بررسی شد خود یک زیرمسأله ساده از \en{IP} است که ماتریس ضرایب آن ساختار ویژه‌ای دارد.
\end{itemize}
\item \textbf{کلاس $P$ در برابر $NP$}: اگر برای یک \en{IP} مشخص الگوریتم چندجمله‌ای یافت شود، صرفاً نشان می‌دهد که آن محدودیت خاص از $IP$ در کلاس $P$ جای دارد، نه آن‌که $IP$ عمومی در $P$ باشد.
\item برای رد \en{NP-hard} بودن کل $IP$، باید ثابت کرد که \textbf{هر} نمونه از برنامه‌ریزی عدد صحیح (از جمله مسائل \en{NP-complete} مانند \en{SAT}, \en{3-SAT}, رنگ‌آمیزی گراف و ... که همگی قابل فرمول‌بندی به صورت $IP$ هستند) در زمان چندجمله‌ای حل می‌شود. این امری است که با شناخت کنونی علوم کامپیوتر معادل با اثبات $P=NP$ خواهد بود.
\end{itemize}

در نتیجه، ارائه یک الگوریتم سریع برای یک مدل $IP$ مشخص، کوچک‌ترین خدشه‌ای به \en{NP-hardness} کلی برنامه‌ریزی عدد صحیح وارد نمی‌کند. \en{NP-hardness} بر اساس وجود \textbf{حداقل یک} نمونه سخت تعریف می‌شود، و پیدا شدن نمونه‌های آسان تأثیری در آن ندارد.

\thickseparator

\end{document}